%%% FIELD phase-sensitive-expected-length-proof
Put
\[
 x=\log(8/\eta),\quad M=mh^d,\quad \ell=hk^{-1/d},
\]
and define
\[
 r_o=n^{-1/(2+d/\gamma)},
 \quad
 r_2=n^{-1/\{1+d/(2\gamma)+d/(2q)\}}.
\tag{1}
\]
The declared domains imply \(d\ge1\), \(q=\alpha+\beta>0\), \(\gamma>0\), and \(x>0\).  They also imply \(0<\epsilon_n\le1\), \(0<\delta_n<1/2\), and \(0<\eta<1/4\), so \(N_{\mathrm{priv}}\) is positive and finite and every division below is legitimate.  Since \(n\ge12\) and \(m=\lfloor n/3\rfloor\),
\[
 {n\over4}\le m\le {n\over3}.
\tag{2}
\]
Because \(s=q/2\), the low-smoothness condition in Definition \(\mathrm{def{:}phase\text{-}rate}\) is
\[
 q<{d\over2\{1+d/(2\gamma)\}}.
\tag{3}
\]
Direct algebra gives
\[
 1+{d\over2\gamma}+{d\over2q}>2+{d\over\gamma}
 \quad\Longleftrightarrow\quad
 q<{d\over2\{1+d/(2\gamma)\}},
\tag{4}
\]
and the two denominators are equal on the boundary.  Hence the piecewise definition of \(r_{\mathrm{np}}\) is equivalently
\[
 r_{\mathrm{np}}=r_o\vee r_2.
\tag{5}
\]
Set
\[
 R=r_{\mathrm{np}}\vee r_{\mathrm{reg}}\vee r_{\mathrm{ho}}
   =r_{\mathrm{up}},
 \quad
 t=1\wedge R.
\tag{6}
\]
Fix a constant \(t_0\) satisfying
\[
 0<t_0<1\wedge r_0^\gamma.
\tag{7}
\]

We first establish, independently of the value of \(t\), that the deterministic minimizer in the statement exists.  Choose \(n_0\ge12\) large enough that, for every \(n\ge n_0\),
\[
 mr_0^d\ge C_S^2x
 \quad\text{and}\quad
 {1\over4}n^{2\gamma/(2\gamma+d)}\ge C_S^2x.
\tag{8}
\]
Such an \(n_0\) exists because \(m=\lfloor n/3\rfloor\) and both left-hand sides diverge with \(n\).  For every \(n\ge n_0\), the bandwidth set is therefore the nonempty compact interval
\[
 I_n=\left[\left\{{C_S^2x\over m}\right\}^{1/d},r_0\right],
\tag{9}
\]
and \((r_0,k_{\mathrm{poly}})\) is feasible, whether or not \(t<t_0\).  Write
\[
 F_n(h,k)=B(h,k)+S(h,k,\eta)
       +G(h,k,\eta,\epsilon_n,\delta_n)
\]
and \(V_n=F_n(r_0,k_{\mathrm{poly}})<\infty\).  The final term of Definition \(\mathrm{def{:}sampling\text{-}envelope}\), together with \(h\le r_0\), yields
\[
 F_n(h,k)
 \ge {C_Sx^2\over m^2h^{2d}}k
 \ge {C_Sx^2\over m^2r_0^{2d}}k
 =:c_nk,
 \quad c_n>0.
\tag{10}
\]
Thus no integer \(k>V_n/c_n\) can minimize \(F_n\).  The remaining candidate integers \(k\ge k_{\mathrm{poly}}\) form a finite nonempty set, and for each such \(k\), continuity of \(F_n(\,\cdot\,,k)\) on the closed bounded interval \(I_n\) gives a minimizing bandwidth.  Minimization over the finite set of candidate integers gives a global minimizer.  Choosing the smallest minimizing \(k\), and then the smallest minimizing \(h\), is a deterministic tie-breaking rule.  Denote the resulting minimizer by \((\widehat h,\widehat k)\).

It remains to bound its objective.  Suppose first that \(t<t_0\).  Then \(t<1\), so (6) gives \(t=R\) and hence
\[
 t\ge r_o,\quad t\ge r_2,
 \quad t\ge r_{\mathrm{reg}},\quad t\ge r_{\mathrm{ho}}.
\tag{11}
\]
For any feasible \((h,k)\), put
\[
 a=M^{-1/2},\quad b={\sqrt{k}\over M}.
\]
Since \(k\ge1\), \(0<\ell\le h\le r_0\le1/4\).  Definitions \(\mathrm{def{:}bias\text{-}envelope}\) and \(\mathrm{def{:}sampling\text{-}envelope}\) therefore imply
\[
 B(h,k)\le C\{h^\gamma+\ell^q\},
\tag{12}
\]
because \(h^\gamma\ell^\alpha\le h^\gamma\), and
\[
 S(h,k,\eta)=C_S\left\{
 \sqrt{x(a^2+b^2)}+xa^2+x^{3/2}ab+x^2b^2
 \right\}.
\tag{13}
\]
Furthermore, the definition of \(N_{\mathrm{priv}}\) and (2) give
\[
 {\sqrt{\log(1/\delta_n)\log(8/\eta)}\over m\epsilon_n}
 ={n\over m}\sqrt{{\log(8/\eta)\over\log(1/\eta)}}
 N_{\mathrm{priv}}^{-1}
 \le C N_{\mathrm{priv}}^{-1}.
\]
Consequently Definition \(\mathrm{def{:}privacy\text{-}envelope}\) gives
\[
 G(h,k,\eta,\epsilon_n,\delta_n)
 \le C N_{\mathrm{priv}}^{-1}\{h^{-d}+kh^{-2d}\}
 =C N_{\mathrm{priv}}^{-1}\{h^{-d}+h^{-d}\ell^{-d}\}.
\tag{14}
\]
Here and below \(C\) denotes a finite constant depending only on the fixed class constants and \(\eta\), and it may change from line to line.

Consider first \(q<\gamma\).  Define the comparison tuning
\[
 h=t^{1/\gamma},\quad \ell_0=t^{1/q},
 \quad K=(h/\ell_0)^d,\quad
 k=\max\{k_{\mathrm{poly}},\lceil K\rceil\}.
\tag{15}
\]
By (7), \(h<r_0\).  Since \(0<t<1\) and \(q<\gamma\), \(K\ge1\), whence
\[
 K\le k\le \max\{2,k_{\mathrm{poly}}\}K,
 \quad
 \max\{2,k_{\mathrm{poly}}\}^{-1/d}\ell_0
 \le \ell\le\ell_0.
\tag{16}
\]
In particular, (12), (15), and (16) show that \(B(h,k)\le Ct\).  Equations (1), (2), and (11) imply
\[
 a=m^{-1/2}t^{-d/(2\gamma)}
 \le2n^{-1/2}t^{-d/(2\gamma)}
 \le2t,
\tag{17}
\]
because the last inequality is equivalent to \(t\ge r_o\).  Similarly, (2), (15), and (16) give
\[
 b\le Cn^{-1}t^{-d/(2\gamma)-d/(2q)}\le Ct,
\tag{18}
\]
where the last inequality is equivalent to \(t\ge r_2\).  Since \(t\le1\), substituting (17)--(18) into (13) yields \(S(h,k,\eta)\le Ct\).

On this branch \(\bar q=q\), so Definition \(\mathrm{def{:}phase\text{-}rate}\) gives
\[
 r_{\mathrm{reg}}=N_{\mathrm{priv}}^{-1/(1+d/\gamma)},
 \quad
 r_{\mathrm{ho}}=N_{\mathrm{priv}}^{-1/(1+d/\gamma+d/q)}.
\tag{19}
\]
Thus the last two inequalities in (11) imply, respectively,
\[
 N_{\mathrm{priv}}^{-1}t^{-d/\gamma}\le t,
 \quad
 N_{\mathrm{priv}}^{-1}t^{-d/\gamma-d/q}\le t.
\tag{20}
\]
Using (14)--(16) in (20) shows that \(G(h,k,\eta,\epsilon_n,\delta_n)\le Ct\).  The comparison tuning (15) is feasible: indeed, (1), (2), and \(t\ge r_o\) give
\[
 mh^d=mt^{d/\gamma}
 \ge {1\over4}n^{2\gamma/(2\gamma+d)}
 \ge C_S^2x,
\tag{21}
\]
where the last inequality follows from the choice (8) of \(n_0\).  We have proved on the \(q<\gamma\) branch that
\[
 F_n(h,k)\le Ct.
\tag{22}
\]
No step imposed an upper bound on \(k\), so this calculation remains valid when \(k>M\).

Now consider \(q\ge\gamma\), still under \(t<t_0\).  Choose the comparison tuning
\[
 h=t^{1/\gamma},\quad k=k_{\mathrm{poly}},
 \quad \ell=hk_{\mathrm{poly}}^{-1/d}.
\tag{23}
\]
Again \(h<r_0\) by (7), and feasibility follows from the same calculation (21).  The integer \(k_{\mathrm{poly}}=(R_*+1)^d\) depends only on the fixed smoothness and dimension parameters.  Hence \(\ell\) is comparable to \(h\); since \(q\ge\gamma\), \(h<1\), and \(\alpha>0\), (12) gives \(B(h,k_{\mathrm{poly}})\le Ct\).  Equation (17) remains valid, while
\[
 b={\sqrt{k_{\mathrm{poly}}}\over M}
   =\sqrt{k_{\mathrm{poly}}}\,a^2\le Ct^2.
\tag{24}
\]
Substitution into (13) gives \(S(h,k_{\mathrm{poly}},\eta)\le Ct\).  Here \(\bar q=\gamma\), and hence
\[
 r_{\mathrm{ho}}
 =N_{\mathrm{priv}}^{-1/(1+2d/\gamma)}
 =N_{\mathrm{priv}}^{-\gamma/(\gamma+2d)}.
\tag{25}
\]
The inequality \(t\ge r_{\mathrm{ho}}\) is equivalent to
\[
 N_{\mathrm{priv}}^{-1}t^{-2d/\gamma}\le t.
\tag{26}
\]
At (23), each term on the right-hand side of (14) is at most a fixed multiple of the left-hand side of (26), because \(0<t<1\) and \(k_{\mathrm{poly}}\) is fixed.  Therefore \(G(h,k_{\mathrm{poly}},\eta,\epsilon_n,\delta_n)\le Ct\), and (22) holds on this branch too.  Equations (19)--(20) identify the \(q\)-branch when \(q<\gamma\), while (25)--(26) identify the private-Gram branch when \(q\ge\gamma\).

By the minimizing property of \((\widehat h,\widehat k)\), (22) implies
\[
 F_n(\widehat h,\widehat k)\le Ct
 \quad\text{whenever }t<t_0.
\tag{27}
\]
Definition \(\mathrm{def{:}private\text{-}confidence\text{-}interval}\) clips its center into \([-1,1]\), and all three envelopes are nonnegative.  The clipped center itself therefore belongs to the intersection defining the interval, so the interval is nonempty on every transcript.  Its length obeys the deterministic bound
\[
 \operatorname{length}\{C_n(\widehat h,\widehat k)\}
 \le 2F_n(\widehat h,\widehat k)\wedge2.
\tag{28}
\]
If \(t<t_0\), equations (27)--(28) bound this length by \(2Ct\).  If \(t\ge t_0\), equation (28) instead gives
\[
 \operatorname{length}\{C_n(\widehat h,\widehat k)\}
 \le2\le {2\over t_0}t.
\tag{29}
\]
The bounds (28)--(29) hold transcript by transcript, uniformly over \(Q\in\mathcal P_{\mathrm{att}}\), and for every deterministic pair of pilot dimensions \(J_\pi,J_\mu\ge1\).  Taking expectation, then the supremum over \(Q\), and recalling from (6) that \(t=1\wedge r_{\mathrm{up}}\), proves the asserted inequality after enlarging the fixed constant to \(C_L\).  The construction uses the fixed lower cutoff \(k\ge k_{\mathrm{poly}}\), imposes no upper relation between \(k\) and \(mh^d\), and makes no claim about a matching full-class frontier.
